\documentclass[12pt,aspectratio=43]{beamer}

%%%%%%%%%%%%%%%%%%%%%%%%%%%%%%%%%%%%%%%%%%%%%%%%%%
%   Paquetes de configuración basica del texto   %
%%%%%%%%%%%%%%%%%%%%%%%%%%%%%%%%%%%%%%%%%%%%%%%%%%
\usepackage[spanish,es-nodecimaldot]{babel}
\usepackage{biblatex}
\usepackage[T1]{fontenc}
\usepackage{listings}
\usepackage{multirow}
\usepackage{hyperref}
\usepackage{hyperxmp}
\usepackage{ifthen}
\usepackage{csquotes}

%%%%%%%%%%%%%%%%%%%%%%%%%%%%%%%%%%%%%%%%%%%%%%%%%%
%                   [listings]                   %
%   Configuración de vizualizacion del código    %
%%%%%%%%%%%%%%%%%%%%%%%%%%%%%%%%%%%%%%%%%%%%%%%%%%
\lstset{
	backgroundcolor=\color{lightgray!10},
	basicstyle=\ttfamily,
	extendedchars=true,
	showspaces=false,
	showstringspaces=false,
	captionpos=b,
	keywordstyle=\bfseries\color{cyan},
	commentstyle=\color{gray},
	stringstyle=\color{orange},
	escapeinside={!>}{<!},
	language=[LaTeX]TeX }

%%%%%%%%%%%%%%%%%%%%%%%%%%%%%%%%%%%%%%%%%%%%%%%%%%
%     Paquetes de configuración graficos del     %
%                   documento                    %
%%%%%%%%%%%%%%%%%%%%%%%%%%%%%%%%%%%%%%%%%%%%%%%%%%
\usepackage{graphicx}
\usepackage{tikz}
\usetikzlibrary{babel}
\usepackage{xcolor}
\usepackage{pdfpages}
\usepackage{animate}

\usepackage{ifxetex}
\ifxetex
	\usepackage[no-math]{fontspec}
	\setmainfont{AncizarSans}[
		Path = ../Fuentes/,
		Extension = .otf,
		BoldFont = *-B,
		ItalicFont = *-I,
		BoldItalicFont = *-BI ]
	\setmonofont{UbuntuMono}[
		Path = ../Fuentes/,
		Extension = .ttf,
		BoldFont = *-B,
		ItalicFont = *-I,
		BoldItalicFont = *-BI ]
	\newcommand{\lmr}{\fontfamily{lmr}\selectfont}
	\newcommand{\lmss}{\fontfamily{lmss}\selectfont}
	\newcommand{\lmtt}{\fontfamily{lmtt}\selectfont}
\fi

%%%%%%%%%%%%%%%%%%%%%%%%%%%%%%%%%%%%%%%%%%%%%%%%%%
%           Configuraciones de Beamer            %
%%%%%%%%%%%%%%%%%%%%%%%%%%%%%%%%%%%%%%%%%%%%%%%%%%
\definecolor{Igreen}{RGB}{148,180,59}
\definecolor{Ired}{RGB}{166,24,49}

\setbeamertemplate{navigation symbols}{}

\setbeamertemplate{itemize items}{\color{Igreen}\raisebox{.45ex}{\rule{.6ex}{.6ex}}}
\setbeamertemplate{enumerate item}{\color{Igreen}\bfseries\arabic{enumi}.}
\setbeamertemplate{enumerate subitem}{\color{Igreen}\bfseries\arabic{enumi}.\arabic{enumii}.}

\usefonttheme{professionalfonts}
\usefonttheme{serif}

\setbeamercolor{frametitle}{fg=Ired,bg=Igreen!50}
\setbeamercolor{alerted text}{fg=Ired}

\setbeamertemplate{bibliography item}{
	\ifboolexpr{ test {\ifentrytype{book}} or test {\ifentrytype{mvbook}}
		or test {\ifentrytype{collection}} or test {\ifentrytype{mvcollection}}
		or test {\ifentrytype{reference}} or test {\ifentrytype{mvreference}} }
	{\setbeamertemplate{bibliography item}[book]}
	{\ifentrytype{online}
		{\setbeamertemplate{bibliography item}[online]}
		{\setbeamertemplate{bibliography item}[article]}}
	\usebeamertemplate{bibliography item}}

\defbibenvironment{bibliography}
	{\list{}
		{\settowidth{\labelwidth}{\usebeamertemplate{bibliography item}}%
			\setlength{\leftmargin}{\labelwidth}%
			\setlength{\labelsep}{\biblabelsep}%
			\addtolength{\leftmargin}{\labelsep}%
			\setlength{\itemsep}{\bibitemsep}%
			\setlength{\parsep}{\bibparsep}}}
	{\endlist}
	{\item}

\makeatletter
\newcommand{\ifratio}[2]{
	\ifthenelse
	{\lengthtest{\beamer@paperwidth=16cm} \AND \lengthtest{\beamer@paperheight=9cm}}
	{#1}
	{#2} }
\makeatother

\setbeamertemplate{title page}{
	\begin{tikzpicture}[overlay,remember picture]
	\path (current page.north west) ++(0.5cm,-0.5cm) node[below right] {\includegraphics[height=1.25cm]{../Escudo_UN}};
	\path (current page.north east) ++(-0.5cm,-0.5cm) node[below left] {\parbox[c][1.25cm][c]{\widthof{\Large \insertsubtitle}}{\Large \insertsubtitle}};
	\end{tikzpicture}

	\vfill
	\begin{center}
		{\Huge\inserttitle}\\
		\bigskip
		\insertauthor
	\end{center}
	\vfill
}

%%%%%%%%%%%%%%%%%%%%%%%%%%%%%%%%%%%%%%%%%%%%%%%%%%
%         Información de la presentación         %
%%%%%%%%%%%%%%%%%%%%%%%%%%%%%%%%%%%%%%%%%%%%%%%%%%

\title{Mas Herramientas de Edición}
\subtitle{Curso Libre de {\lmr\LaTeX}}
\author{Joar Esteban Buitrago Carrillo}
\institute{Universidad Nacional de Colombia}
\date{}
\logo{\includegraphics[height=0.75cm]{../Escudo_UN}}
\hypersetup{
	pdfcopyright={Esta obra está bajo una licencia de Creative Commons Reconocimiento-CompartirIgual 4.0 Internacional.},
	pdflicenseurl={https://creativecommons.org/licenses/by-sa/4.0/deed.es}
}
%\addbibresource{biblio.bib}

\begin{document}
\begin{frame}[plain]
\titlepage
\end{frame}

{
\setbeamercolor{background canvas}{bg=}
\begin{frame}
\ifratio
	{\includepdf[pages=2]{../Licencia_Diapositiva.pdf}}
	{\includepdf[pages=1]{../Licencia_Diapositiva.pdf}}
\end{frame}
}

\section{Referencias Cruzadas}
\begin{frame}{Referencias Cruzadas}{}
Como en la mayoría de editores de texto es necesario el uso de referencias cruzadas para citar una gráfica, sección o cualquier parte del texto.\pause\\[1em]

En {\lmr\LaTeX} esto funciona por medio de etiquetas que se incluyen en el código fuente. De esta forma se vuelve algo muy sencillo aprenderse una etiqueta y no un número.
\end{frame}

\begin{frame}[fragile]{Referencias Cruzadas}{}
\begin{itemize}[<+->]
	\item Para colocar una etiqueta se utiliza el comando \lstinline|\label{<etiqueta>}|
	\item Para citar el número asociado a una etiqueta se utiliza el comando \lstinline|\ref{<etiqueta>}|
	\item Para citar la página donde se encuentra una etiqueta se utiliza el comando \lstinline|\pageref{<etiqueta>}|
\end{itemize}
\end{frame}

\begin{frame}[fragile]{Referencias Cruzadas}{}
Así, por ejemplo:

\begin{columns}
	\begin{column}{0.5\linewidth}
		{\small
\begin{lstlisting}
\section{Sección a}
\label{sec:a}
Texto de la sección...

Podemos ver en la
seccion \ref{sec:a}
\end{lstlisting} }
	\end{column}
	\vline
	\hspace{1em}
	\begin{column}{0.5\linewidth}
		{\lmr{\small
				{\large\bfseries 5. Sección A}\\[2.3ex]
				Texto de la sección...\\[2em]
				Podemos ver en la sección 5 }}
	\end{column}
\end{columns}
{\color{Ired} Código}\hfill{\color{Ired} Vista}
\end{frame}

\section{Texto Cargado}
\begin{frame}{Texto Cargado}{}
{\lmr\LaTeX} es capaz de elegir la justificación de un texto, para ello, podemos hacer uso de comandos globales, así como de entornos para cambiar la justificación de un texto.\pause\\[1em]

Para hacer este proceso mas facil vamos a utilizar el paquete \alert{\texttt{ragged2e}}.
\end{frame}

\subsection{Cargado a izquierda}
\begin{frame}{Cargado a izquierda}{Texto Cargado}
Lo más normal es cargar un texto a la izquierda, es lo que solemos hacer al escribir en un papel.\pause\\[1em]

Para hacer esto en {\lmr\LaTeX} se utiliza el comando \texttt{\textbackslash raggedright} («Roto a derecha») o el entorno \texttt{flushleft}.
\end{frame}

\begin{frame}[fragile]{Cargado a izquierda}{Texto Cargado}
Así, por ejemplo:

\begin{columns}
	\begin{column}{0.5\linewidth}
		{\small
\begin{lstlisting}
Este texto está
justificado como se
puede ver en el margen
derecho

\raggedright
Este texto está cargado
a izquierda
\end{lstlisting} }
	\end{column}
	\vline
	\hspace{1em}
	\begin{column}{0.5\linewidth}
		{\lmr{\small
				\parbox{\linewidth}{Este texto está justificado como se puede ver en el margen derecho}\\[2em]
				Este texto está cargado a izquierda }}
	\end{column}
\end{columns}
{\color{Ired} Código}\hfill{\color{Ired} Vista}
\end{frame}

\begin{frame}[fragile]{Cargado a izquierda}{Texto Cargado}
\begin{columns}
\begin{column}{0.5\linewidth}
	{\small
\begin{lstlisting}
Este texto está
justificado como se
puede ver en el margen
derecho

\begin{flushleft}
  Este texto está
  cargado a izquierda
\end{flushleft}
\end{lstlisting} }
\end{column}
\vline
\hspace{1em}
\begin{column}{0.5\linewidth}
	{\lmr{\small
			\parbox{\linewidth}{Este texto está justificado como se puede ver en el margen derecho}\\[2em]
			Este texto está cargado a izquierda }}
\end{column}
\end{columns}
{\color{Ired} Código}\hfill{\color{Ired} Vista}
\end{frame}

\subsection{Cargado a derecha}
\begin{frame}{Cargado a derecha}{Texto Cargado}
Ahora, para cargar un texto a derecha en {\lmr\LaTeX} se utiliza el comando \texttt{\textbackslash raggedleft} («Roto a izquierda») o el entorno \texttt{flushright}.
\end{frame}

\begin{frame}[fragile]{Cargado a derecha}{Texto Cargado}
Así, por ejemplo:

\begin{columns}
\begin{column}{0.5\linewidth}
	{\small
\begin{lstlisting}
Este texto está
justificado como se
puede ver en el margen
derecho

\raggedleft
Este texto está cargado
a derecha
\end{lstlisting} }
\end{column}
\vline
\hspace{1em}
\begin{column}{0.5\linewidth}
	{\lmr{\small
			\parbox{\linewidth}{Este texto está justificado como se puede ver en el margen derecho}\\[2em]
			\null\hfill Este texto está cargado a\\\vspace*{-3ex}
			\null\hfill derecha\null }}
\end{column}
\end{columns}
{\color{Ired} Código}\hfill{\color{Ired} Vista}
\end{frame}

\begin{frame}[fragile]{Cargado a derecha}{Texto Cargado}
\begin{columns}
\begin{column}{0.5\linewidth}
	{\small
\begin{lstlisting}
Este texto está
justificado como se
puede ver en el margen
derecho

\begin{flushright}
  Este texto está
  cargado a derecha
\end{flushright}
\end{lstlisting} }
\end{column}
\vline
\hspace{1em}
\begin{column}{0.5\linewidth}
	{\lmr{\small
			\parbox{\linewidth}{Este texto está justificado como se puede ver en el margen derecho}\\[2em]
			\null\hfill Este texto está cargado a\\\vspace*{-3ex}
			\null\hfill derecha\null }}
\end{column}
\end{columns}
{\color{Ired} Código}\hfill{\color{Ired} Vista}
\end{frame}

\subsection{Texto centrado}
\begin{frame}{Texto centrado}{Texto Cargado}
Para centrar un texto no es muy diferente, se utiliza el comando \texttt{\textbackslash centering} o el entorno \texttt{center}.
\end{frame}

\begin{frame}[fragile]{Texto centrado}{Texto Cargado}
Así, por ejemplo:

\begin{columns}
\begin{column}{0.5\linewidth}
	{\small
\begin{lstlisting}
Este texto está
justificado como se
puede ver en el margen
derecho

\centering
Este texto está
centrado
\end{lstlisting} }
\end{column}
\vline
\hspace{1em}
\begin{column}{0.5\linewidth}
	{\lmr{\small
			\parbox{\linewidth}{Este texto está justificado como se puede ver en el margen derecho}\\[2em]
			\null\hfill Este texto está centrado\hfill\null}}
\end{column}
\end{columns}
{\color{Ired} Código}\hfill{\color{Ired} Vista}
\end{frame}

\begin{frame}[fragile]{Cargado a izquierda}{Texto Cargado}
\begin{columns}
\begin{column}{0.5\linewidth}
	{\small
\begin{lstlisting}
Este texto está
justificado como se
puede ver en el margen
derecho

\begin{center}
  Este texto está
  centrado
\end{center}
\end{lstlisting} }
\end{column}
\vline
\hspace{1em}
\begin{column}{0.5\linewidth}
{\lmr{\small
		\parbox{\linewidth}{Este texto está justificado como se puede ver en el margen derecho}\\[2em]
		\null\hfill Este texto está centrado\hfill\null }}
\end{column}
\end{columns}
{\color{Ired} Código}\hfill{\color{Ired} Vista}
\end{frame}

\section{Contadores}
\begin{frame}{Contadores}{}
Los contadores son uno de los elementos más importantes de {\lmr\LaTeX}, estos controlan el número de la página, sección, el capitulo, etc. de lo que estés trabajando.\pause\\[1em]

Los contadores pueden ser modificados y creados.
\end{frame}

\begin{frame}{Contadores}{}
Para acceder a el valor de un contador, se utiliza la orden \texttt{\textbackslash the\textit{contador}}.\pause\\[1em]

Entonces, por ejemplo, para acceder al número de la página utilizamos el comando \texttt{\textbackslash thepage}, dando como resultado \thepage.\pause\\[1em]

Cada uno de los elementos de {\lmr\LaTeX} suele tener un contador asociado con el mismo nombre.
\end{frame}

\begin{frame}[fragile]{Contadores}{}
Para la modificación y creación de contadores existen 4 comandos básicos.

\begin{itemize}[<+->]
	\item \lstinline|\setcounter{<contador>}{<número>}| Con lo que el contador toma el valor del número, que es un entero positivo o negativo.
	\item \lstinline|\addtocounter{<contador>}{<número>}| Con lo que el contador incrementa en la cantidad indicada por el número.
\end{itemize}
\end{frame}

\begin{frame}[fragile]{Contadores}{}
Para la modificación y creación de contadores existen 4 comandos básicos.

\begin{itemize}[<+->]
	\item \lstinline|\stepcounter{<contador>}{<número>}| o \lstinline|\refstepcounter{<contador>}{<número>}| Con lo cual el contador incrementa en uno. Con la versión \texttt{ref} es posible usar referencias cruzadas.
	\item \lstinline|\newcounter{<contador>}[<suborden>]| Con lo cual se crea el contador, y se le asigna, de ser necesario, una subordinación a otro contador.
\end{itemize}
\end{frame}

\begin{frame}{Contadores}{}
Los contadores que {\lmr\LaTeX} define por defecto son

\begin{center}
	\begin{tabular}{llll}
		part & paragraph & figure & enumi \\
		chapter & subparagraph & table & enumii \\
		section & page & footnote & enumiii \\
		subsection & equation & mpfootnote & enumiv \\
		subsubsection &  &  &  \\
	\end{tabular}
\end{center}
\end{frame}

\begin{frame}[fragile]{Contadores}
Para cambiar el tipo de numeración que maneja un contador, {\lmr\LaTeX} posee las ordenes

\begin{itemize}
	\item \lstinline|\arabic{<contador>}|
	\item \lstinline|\roman{<contador>}|
	\item \lstinline|\Roman{<contador>}|
	\item \lstinline|\alph{<contador>}|
	\item \lstinline|\Alph{<contador>}|
\end{itemize}\pause

Entonces, para cambiar el tipo de numeración de un contador escribimos en el código por ejemplo \lstinline|\renewcommand{\the<contador>}{\Roman{<contador>}}|
\end{frame}

\section{Comandos y Entornos Personalizados}
\begin{frame}{Comandos y Entornos Personalizados}{}
Cuando utilizamos mucho un comando, o una serie de comandos, y es tedioso escribirlos, porque son muy largos, lo mejor es crear un nuevo comando, con todo lo que necesitamos.\\[1em]

Esto también aplica con entornos y grupos.
\end{frame}

\subsection{Comandos Personalizados}
\begin{frame}[fragile]{Comandos Personalizados}{Comandos y Entornos Personalizados}
Para crear un comando personalizado utilizamos las ordenes

\begin{itemize}[<+->]
	\item \lstinline|\newcommand{\<comando>}|\\
	\texttt{[<\# Argumentos>][<Argumento por defecto>]\{<Definición>\}}
	\item \lstinline|\renewcommand{\<comando>}|\\
	\texttt{[<\# Argumentos>][<Argumento por defecto>]\{<Definición>\}}
	\item \lstinline|\providecommand{\<comando>}|\\
	\texttt{[<\# Argumentos>][<Argumento por defecto>]\{<Definición>\}}
\end{itemize}
\end{frame}

\begin{frame}[fragile]{Comandos Personalizados}{Comandos y Entornos Personalizados}
Por ejemplo:\\[1em]

\lstinline|$\mathbb{R}^n$| produce $\mathbb{R}^n$ y si utilizamos este comando muchas veces en el documento lo mejor seria crear un nuevo comando que simplifique esto, así:\\[1em]

\begin{lstlisting}
\newcommand{\R}[1][2]{\mathbb{R}^#1}
\end{lstlisting}\vspace{1em}

Con lo cual solo necesitamos escribir \lstinline|$\R[2]$| para obtener  $\mathbb{R}^2$
\end{frame}

\subsection{Entornos Personalizados}
\begin{frame}[fragile]{Entornos Personalizados}{Comandos y Entornos Personalizados}
Para crear un entorno personalizado utilizamos las ordenes

\begin{itemize}[<+->]
	\item \lstinline|\newenvironment{<nombre>}|\\
	\texttt{[<\# Argumentos>][<Argumento por defecto>]\{<Parte inicial>\}\{Parte final\}}
	\item \lstinline|\renewenvironment{<nombre>}|\\
	\texttt{[<\# Argumentos>][<Argumento por defecto>]\{<Parte inicial>\}\{Parte final\}}
\end{itemize}
\end{frame}

\begin{frame}[fragile]{Entornos Personalizados}{Comandos y Entornos Personalizados}
Por ejemplo:\\[1em]

\begin{lstlisting}
\newenvironment{miTabla}[][]{
  \begin{tabular}{|c|c|c|}
}{
  \end{tabular}
}
\end{lstlisting}\vspace{1em}

Con lo cual solo necesitamos utilizar el entorno \texttt{miTabla} para crear una tabla con 3 columnas centradas.
\end{frame}

\section{El paquete \texttt{showkeys}}
\end{document}
